% Part: computability
% Chapter: computability-theory
% Section: equiv-ce-defs

\documentclass[../../../include/open-logic-section]{subfiles}

\begin{document}

\olfileid{cmp}{thy}{eqc}

\olsection[د c.e. سټونو تعريفونه]{په محاسبوي ډول د شمېر وړ سټونو
  معادل تعريفونه}
\textbf{د ژباړې د نسخې سمون (OLCMP-024):} د سرچينې په اوږده سرليک
کښې د تعريفونو د انګرېزي نوم املا دوه توري زيات لرل، او لاندې تشريح کښې د
``هغه ننوتونه چې محاسبه پرې درېږي'' پر ځاے دوه سربلونه سره نښتي وو۔
په پښتو متن کښې دواړه څرګندې سمونيزې تېروتنې سمې شوې دي؛ هېڅ فني
دعوه يا نښه نۀ ده بدله شوې۔


لاندې قضيه د دې خبرې څو مهم معادل بيانونه راکوي چې يو سټ په محاسبوي
ډول د شمېر وړ وي۔

\begin{thm}
\ollabel{thm:ce-equiv}
فرض کړئ $S$ د طبيعي عددونو يو سټ دے۔ بيا لاندې خبرې سره معادلې دي:
\begin{enumerate}
\item\ollabel{case:ce} $S$ په محاسبوي ډول د شمېر وړ دے۔
\item\ollabel{case:ran-pc} $S$ د يوې \emph{جزوي} محاسبه کېدونکې تابع
د قيمتونو سټ دے۔
\item\ollabel{case:ran-prim} $S$ تش دے يا د يوې بنسټيزې بازګشتي تابع
د قيمتونو سټ دے۔
\item\ollabel{case:ce-domain} $S$ د يوې جزوي محاسبه کېدونکې تابع
\emph{د تعريف ساحه} ده۔
\end{enumerate}
\end{thm}

\begin{explain}
لومړۍ درې فقرې وايي چې هر ناتش په محاسبوي ډول د شمېر وړ سټ په معادل
ډول د يوې محاسبه کېدونکې تابع، جزوي محاسبه کېدونکې تابع، يا بنسټيزې
بازګشتي تابع په وسيله شمېرل کېدے شي۔ څلورمه فقره وايي چې که~$S$ په
محاسبوي ډول د شمېر وړ وي، نو د کوم شاخص~$e$ دپاره
\[
S = \Setabs{x}{\cfind{e}(x) \fdefined}.
\]
په بله وينا، $S$ د هغو ننوتونو سټ دے چې د~$\cfind{e}$ محاسبه پرې
درېږي۔ له همدې کبله په محاسبوي ډول د شمېر وړ سټونه کله ناکله
\emph{نيمه د پرېکړې وړ} بلل کېږي: که يو عدد په سټ کښې وي، بالاخره
``هو'' تر لاسه کوئ؛ خو که پکښې نۀ وي، هېڅکله ``نه'' نۀ تر لاسه
کوئ!{}
\end{explain}

\begin{proof}
څرنګه چې هره بنسټيزه بازګشتي تابع محاسبه کېدونکې ده او هره محاسبه
کېدونکې تابع جزوي محاسبه کېدونکې ده، \olref{case:ran-prim} پر
\olref{case:ce} دلالت کوي او \olref{case:ce} پر
\olref{case:ran-pc} دلالت کوي۔ (پام وکړئ چې که~$S$ تش وي، نو~$S$ د
هغې جزوي محاسبه کېدونکې تابع د قيمتونو سټ دے چې هېڅ ځاے تعريف شوې
نۀ ده۔) که وښيو چې \olref{case:ran-pc} پر
\olref{case:ran-prim} دلالت کوي، لومړۍ درې فقرې مو سره معادلې ښودلې
دي۔

نو فرض کړئ $S$ د جزوي محاسبه کېدونکې تابع~$\cfind{e}$ د قيمتونو سټ
دے۔ که~$S$ تش وي، کار بشپړ دے۔ که نه، $a$ د~$S$ هر يو غړے وي۔ د
کليني د نورمال شکل د قضيې له مخې ليکلے شو:
\[
\cfind{e}(x) = U(\umin{s}{T(e, x, s)}).
\]
په ځانګړي ډول، $\cfind{e}(x) \fdefined$ او $= y$ هله او يوازې هله
چې داسې~$s$ وي چې $T(e, x, s)$ او $U(s) = y$ وي۔ تابع~$f(z)$ داسې
تعريف کړئ:
\[
f(z) = \begin{cases}
  U((z)_1) & \text{که $T(e, (z)_0, (z)_1)$ وي} \\
  a        & \text{که داسې نۀ وي۔}
\end{cases}
\]
بيا~$f$ بنسټيزه بازګشتي ده، ځکه~$T$ او~$U$ داسې دي۔
\iftag{TMs}{د ټيورينګ ماشينونو په ژبه، که~$z$ داسې جوړه
  $\tuple{(z)_0, (z)_1}$ کوډوي چې $(z)_1$ د ماشين~$M_e$ پر
  ننوت~$(z)_0$ يوه درېدلې محاسبه وي، نو~$f$ د محاسبې وت راګرځوي؛
  که نه، $a$ راګرځوي۔}

بايد وښيو چې~$S$ د~$f$ د قيمتونو سټ دے؛ يعنې د هر طبيعي عدد~$y$
دپاره، $y \in S$ هله او يوازې هله چې د~$f$ په قيمتونو کښې وي۔ په
مخکښ لوري کښې فرض کړئ $y \in S$۔ بيا~$y$ د~$\cfind{e}$ په قيمتونو
کښې دے، نو د ځينو~$x$ او~$s$ دپاره $T(e,x,s)$ پوره کېږي او
$U(s) = y$ ده۔ خو بيا $y = f(\tuple{x,s})$۔ بل پلو، فرض کړئ~$y$ د
~$f$ په قيمتونو کښې دے۔ بيا يا $y = a$، يا د کوم~$z$ دپاره
$T(e,(z)_0,(z)_1)$ او $U((z)_1) = y$ دي۔ څرنګه چې په وروستي حالت
کښې $\cfind{e}((z)_0) \fdefined = y$ ده، په دواړو حالونو کښې~$y$
په~$S$ کښې دے۔
\textbf{د ژباړې د نسخې سمون (OLCMP-025):} سرچينې په وروستۍ نتيجه
کښې يو ناتړلے متغير کارولے و، که څه هم د شاهد کوډ لومړۍ برخه د
محاسبې ننوت ټاکي۔ دلته د نورمال شکل له مخکنۍ جملې سره سم د شاهد له
کوډ څخه اخيستل شوے ننوت ليکل شوے دے۔

(نښه $\cfind{e}(x) \fdefined = y$ دا مانا لري چې ``$\cfind{e}(x)$
تعريف شوې او له~$y$ سره برابره ده۔'' موږ په همدې اسانتيا
$\cfind{e}(x) = y$ هم کارولے شو؛ خو زيات غشے کله ناکله دا رايادوي
چې له يوې جزوي تابع سره کار لرو۔)

د \olref{thm:ce-equiv} د ثبوت د بشپړولو دپاره دومره بسنه کوي چې
وښيو \olref{case:ce} او~\olref{case:ce-domain} سره معادل دي۔ لومړے
ښيو چې \olref{case:ce} پر~\olref{case:ce-domain} دلالت کوي۔ که ورکړل
شوے سټ تش وي، د هغې جزوي محاسبه کېدونکې تابع د تعريف ساحه ده چې هېڅ
ځاے تعريف شوې نۀ ده۔ اوس د ناتش حالت دپاره فرض کړئ~$S$ د يوې محاسبه
کېدونکې تابع~$f$ د قيمتونو سټ دے، يعنې
\[
S = \Setabs{y}{\text{د کوم $x$ دپاره، } f(x) = y}.
\]
وټاکئ:
\[
g(y) = \umin{x}{(f(x) = y)}.
\]
بيا~$g$ جزوي محاسبه کېدونکې تابع ده، او $g(y)$ هله او يوازې هله
تعريف شوې ده چې د کوم~$x$ دپاره $f(x) = y$ وي۔ په بله وينا، د~$g$
د تعريف ساحه د~$f$ د قيمتونو سټ دے۔
\iftag{TMs}{د ټيورينګ ماشينونو په ژبه: که ټيورينګ ماشين~$F$ د~$S$
غړي شمېري، نو ټيورينګ ماشين~$G$ داسې وټاکئ چې په~$S$ کښې د ورکړل
شوي غړي د موندلو دپاره د~$F$ په وتونو کښې پلټنه وکړي؛ که ئې ومومي
ودرېږي، او که ئې نۀ مومي د تل دپاره پلټنې ته دوام ورکړي۔}{}
\textbf{د ژباړې د نسخې سمون (OLCMP-026):} سرچينې په دې دلالت کښې
سمدستي فرض کړې وه چې ورکړل شوے سټ د يوې هرځاے تعريف شوې محاسبه
کېدونکې تابع د قيمتونو سټ دے؛ دا د تش سټ دپاره ناشوني دي، که څه هم
تعريف تش سټ په څرګند ډول c.e. مني۔ دلته تش حالت لومړے د هېڅ ځاے
تعريف شوې جزوي تابع په وسيله حل شوے، او د تابع د قيمتونو استدلال
ناتش حالت ته کارول شوے دے۔

په پای کښې، د دې د ښودلو دپاره چې \olref{case:ce-domain} پر
\olref{case:ce} دلالت کوي، فرض کړئ~$S$ د جزوي محاسبه کېدونکې
تابع~$\cfind{e}$ د تعريف ساحه ده، يعنې
\[
S = \Setabs{x}{\cfind{e}(x) \fdefined}.
\]
که~$S$ تش وي، کار بشپړ دے؛ که نه، $a$ د~$S$ هر يو غړے وي۔ تابع~$f$
داسې تعريف کړئ:
\[
f(z) = \begin{cases}
(z)_0 & \text{که $T(e,(z)_0,(z)_1)$ وي} \\
a & \text{که داسې نۀ وي۔}
\end{cases}
\]
بيا، لکه پورته، عدد~$x$ هله او يوازې هله د~$f$ په قيمتونو کښې دے
چې $\cfind{e}(x) \fdefined$ وي، يعنې هله او يوازې هله چې $x \in S$
وي۔ \iftag{TMs}{د ټيورينګ ماشينونو په ژبه: که ماشين~$M_e$ د~$S$
نيمه پرېکړه کوي، د~$S$ غړي د ټيورينګ ماشينونو د ټولو ممکنه محاسبو
په کتلو او د درېدلو محاسبو د اړوندو ننوتونو په راګرځولو سره وشمېرئ۔}{}
\end{proof}

د \olref{thm:ce-equiv} فقره~\olref{case:ce-domain} موږ ته په
محاسبوي ډول د شمېر وړ سټونو د شمېرلو يوه اسانه لار راکوي: د هر~$e$
دپاره~$W_e$ د~$\cfind{e}$ د تعريف ساحه وټاکئ، يعنې
\[
W_e = \Setabs{x}{\cfind{e}(x) \fdefined}.
\]
بيا که~$A$ هر په محاسبوي ډول د شمېر وړ سټ وي، د کوم~$e$ دپاره
$A = W_e$ وي۔

لاندې قضيه په محاسبوي ډول د شمېر وړ سټونو يوه بله ځانګړنه راکوي۔

\begin{thm}
\ollabel{thm:exists-char}
سټ~$S$ هله او يوازې هله په محاسبوي ډول د شمېر وړ دے چې داسې محاسبه
کېدونکې اړيکه~$R(x,y)$ وي چې
\[
S = \Setabs{ x }{ \lexists[y][R(x,y)] }.
\]
\end{thm}

\begin{proof}
په مخکښ لوري کښې فرض کړئ~$S$ په محاسبوي ډول د شمېر وړ دے۔ نو د کوم
~$e$ دپاره $S = W_e$۔ د~$e$ د همدې قيمت دپاره~$S$ داسې ليکلے شو:
\[
S = \Setabs{ x }{ \lexists[y][T(e, x, y)] }.
\]
په بېرته لوري کښې فرض کړئ $S = \Setabs{ x }{
  \lexists[y][R(x, y)] }$۔ تابع~$f$ داسې تعريف کړئ:
\[
f(x) \simeq \umin{y}{R(x, y)}.
\]
بيا~$f$ جزوي محاسبه کېدونکې ده، او~$S$ د~$f$ د تعريف ساحه ده۔
\end{proof}


\end{document}
